\section{The Two Closest Per-User Training Methods}
\label{sec:closest-per-user}

\subsection{What counts as a close match}

Our unit of personalization is one user.  A deployment run owns that user's
adapter, optimizer state, memory, and checkpoint history.  The initial policy
is trained from the user's thoughts, conversations, corrections, and task
records.  Later interactions update the same policy.  A shared backbone may
be trained separately, but the target user's raw history does not train
another user's deployed branch.

Under this definition, the two closest papers are OPPU
\citep{tan2024oppu} and SPRInG \citep{kim2026spring}.  OPPU is the closest
offline initializer: it trains one adapter from one user's history.  SPRInG
is the closest continual method: it carries that user's adapter forward and
updates it from chronological interaction batches.  Together they bracket
our offline-to-online workflow.

\begin{table}[t]
  \centering
  \caption{Why OPPU and SPRInG are the closest matches.}
  \label{tab:closest-match}
  \begin{tabularx}{\textwidth}{@{}L{0.18\textwidth}YYY@{}}
    \toprule
    Property & OPPU & SPRInG & Our setting \\
    \midrule
    Policy owner &
    One PEFT module per user &
    One evolving LoRA adapter per user &
    One user-owned policy branch \\
    Historical input &
    Paired behaviors or plain text &
    Ordered query--response records &
    Thoughts, conversations, task trajectories, and outcomes \\
    Offline update &
    Cross-entropy or next-token training &
    Pretrained period-0 adapter &
    User-history initialization \\
    Later update &
    None &
    Selective SFT on each new period &
    Preference or outcome-driven online update \\
    Persistent output &
    Personal adapter &
    Adapter lineage and bounded memory &
    Adapter, optimizer state, memory, and audit receipt \\
    \bottomrule
  \end{tabularx}
\end{table}

We exclude methods that train one policy across many users and personalize
only through a profile or prompt.  They do not produce the user-owned policy
studied here.

\subsection{OPPU: one history, one adapter}

\paragraph{Model and data.}
OPPU freezes a task-adapted base model with parameters $\theta$ and assigns
user $u$ a private PEFT module $\Delta\theta_u$.  The deployed model is
\begin{equation}
  M_u = M_{\theta\oplus\Delta\theta_u}.
  \label{eq:oppu-model}
\end{equation}
The user's history $H_u^{<t}$ contains only records observed before the
current query.  A record is either a supervised pair $(x_i,y_i)$ or a plain
text sequence $x_i$.  The first form covers prior requests and accepted
responses.  The second can hold notes, conversation text, or other
user-authored material.

\paragraph{Workflow.}
\begin{enumerate}
  \item Adapt a shared base model to the task using users outside the target
  cohort, then merge the task LoRA into the base.
  \item Freeze the shared base.
  \item Initialize one LoRA module $\Delta\theta_u$ for the target user.
  \item Train it only on $H_u$.  When retrieval is enabled, each example may
  retrieve only records that preceded that example.
  \item Save $\Delta\theta_u$.  At inference, load the frozen base and this
  user's adapter, then answer the new query.
\end{enumerate}

\paragraph{Optimization.}
For paired history, OPPU minimizes conditional cross-entropy:
\begin{equation}
  \loss_{\mathrm{pair}}(\Delta\theta_u)
  =
  -\sum_{(x_i,y_i)\in H_u}
   \sum_{j=1}^{|y_i|}
   \log p_{\theta,\Delta\theta_u}
   \left(y_{i,j}\mid \varphi(x_i),y_{i,<j}\right).
  \label{eq:oppu-paired}
\end{equation}
When history contains plain text only, the target is the next token:
\begin{equation}
  \loss_{\mathrm{text}}(\Delta\theta_u)
  =
  -\sum_{x_i\in H_u}
   \sum_{j=1}^{|x_i|}
   \log p_{\theta,\Delta\theta_u}
   \left(x_{i,j}\mid x_{i,<j}\right).
  \label{eq:oppu-text}
\end{equation}
The paper uses rank-8 LoRA on the query and value projections, learning rate
$10^{-5}$, and two or three fixed epochs.  It does not report an online
update, replay rule, rollback rule, or optimizer-state lifecycle.

\begin{algorithm}[t]
  \caption{OPPU for one user}
  \label{alg:oppu-one-user}
  \KwIn{Frozen task model $M_\theta$; chronological user history $H_u$;
  optional retriever $R$}
  \KwOut{Personal adapter $\Delta\theta_u$}
  initialize $\Delta\theta_u$\;
  \ForEach{record $h_i\in H_u$ in chronological order}{
    construct its target from $(x_i,y_i)$ or right-shifted text $x_i$\;
    if retrieval is enabled, retrieve only from $H_u^{<i}$\;
    accumulate the loss in \cref{eq:oppu-paired} or \cref{eq:oppu-text}\;
  }
  update only $\Delta\theta_u$ for a fixed number of epochs\;
  save $\Delta\theta_u$ with its user and history-version identifiers\;
\end{algorithm}

\paragraph{Inputs and outputs.}
Training consumes a frozen base model, one user identifier, and that user's
historical records.  It outputs one loadable adapter.  Inference consumes the
adapter and a new query; optional retrieval or a generated profile may also
enter the prompt.  The paper evaluates Llama-2-7B on seven LaMP tasks using
the 100 users with the longest histories.

\paragraph{What transfers.}
OPPU gives us the correct ownership boundary and a practical offline
initializer.  It does not learn from task outcomes, pairwise preferences, or
live feedback.  It also treats each record as text rather than as a replayable
agent trajectory.  Our offline stage can use its per-user adapter design, but
our record schema must retain actions, tools, outcomes, timestamps, and the
generating checkpoint.

\subsection{SPRInG: one adapter through time}

\paragraph{Model and data.}
SPRInG starts from the same ownership model.  User $u$ has a frozen backbone
$M_\theta$, an adapter $\phi_u^{(t)}$, and a bounded memory
$B_u^{(t)}$.  New interactions arrive in periods:
\begin{equation}
  D_u^{(t)}=\{(q_i,a_i)\}_{i=1}^{n_t}.
  \label{eq:spring-stream}
\end{equation}
Period 0 supplies the initial adapter.  Each later period updates that same
adapter.  The paper uses text-generation records, not tool trajectories or
explicit rewards.

\paragraph{Workflow.}
\begin{enumerate}
  \item Score every new $(q,a)$ under the frozen base and the current
  user-adapted model.
  \item Keep the top 30\% by drift score.
  \item Copy $\phi_u^{(t-1)}$ and fine-tune the copy on the selected records.
  \item Merge selected records with the old memory, score them again under
  the updated adapter, and retain the hardest 50.
  \item For a new query, retrieve two memory records.  Use them only when
  their similarity passes a fixed gate.
  \item Decode by interpolating adapter-only and retrieval-conditioned token
  probabilities.
\end{enumerate}

\paragraph{Selection and update.}
Let
\begin{align}
  p_{\mathrm{base}}(q,a)
  &=P(a\mid q;M_\theta),\\
  p_{\mathrm{adap}}^{(t-1)}(q,a)
  &=P(a\mid q;M_{\theta,\phi_u^{(t-1)}}).
\end{align}
SPRInG scores a record by
\begin{equation}
  S_t(q,a)
  =
  -\log\frac{p_{\mathrm{adap}}^{(t-1)}(q,a)}
                  {p_{\mathrm{base}}(q,a)}
  +\alpha\log p_{\mathrm{base}}(q,a).
  \label{eq:spring-score}
\end{equation}
The likelihood ratio favors behavior not yet captured by the adapter.  The
base likelihood suppresses malformed or low-quality text.  If $\tau_p$ is
the period's percentile threshold,
\begin{equation}
  D_{\mathrm{drift}}^{(t)}
  =
  \{(q,a)\in D_u^{(t)}:S_t(q,a)\ge\tau_p\}.
  \label{eq:spring-drift}
\end{equation}
The adapter then takes supervised next-token updates on
$D_{\mathrm{drift}}^{(t)}$ only:
\begin{equation}
  \phi_u^{(t)}
  =
  \operatorname{SFT}
  \left(\phi_u^{(t-1)},D_{\mathrm{drift}}^{(t)}\right).
  \label{eq:spring-update}
\end{equation}

\paragraph{Residual memory and inference.}
After the update, SPRInG rescans
\begin{equation}
  C_u^{(t)}
  =
  D_{\mathrm{drift}}^{(t)}\cup B_u^{(t-1)}
\end{equation}
and recomputes \cref{eq:spring-score} with $\phi_u^{(t)}$.  It keeps the
$N_{\max}$ highest-scoring records:
\begin{equation}
  B_u^{(t)}
  =
  \operatorname{Top}_{N_{\max}}
  \{(q,a)\in C_u^{(t)};S'_t(q,a)\}.
  \label{eq:spring-buffer}
\end{equation}
For a test query, the same adapted model produces an adapter-only
distribution $p_{\mathrm{par}}$ and, when memory passes the relevance gate, a
retrieval-conditioned distribution $p_{\mathrm{ret}}$.  The final
distribution is
\begin{equation}
  p_{\mathrm{fin}}
  =
  \lambda p_{\mathrm{ret}}+(1-\lambda)p_{\mathrm{par}}.
  \label{eq:spring-interpolation}
\end{equation}
Memory affects inference only.  It is not replayed into later optimizer
steps.

\begin{algorithm}[t]
  \caption{SPRInG for one user}
  \label{alg:spring-one-user}
  \KwIn{Frozen $M_\theta$; initial $\phi_u^{(0)}$; stream
  $\{D_u^{(t)}\}_{t\ge1}$; percentile $p$; budget $N_{\max}$}
  \KwOut{Adapter lineage $\{\phi_u^{(t)}\}$ and memories
  $\{B_u^{(t)}\}$}
  $B_u^{(0)}\leftarrow\varnothing$\;
  \For{$t=1,2,\ldots$}{
    score $D_u^{(t)}$ with \cref{eq:spring-score}\;
    retain $D_{\mathrm{drift}}^{(t)}$ with \cref{eq:spring-drift}\;
    $\phi_u^{(t)}\leftarrow
      \operatorname{SFT}(\phi_u^{(t-1)},D_{\mathrm{drift}}^{(t)})$\;
    $C_u^{(t)}\leftarrow D_{\mathrm{drift}}^{(t)}\cup B_u^{(t-1)}$\;
    rescore $C_u^{(t)}$ under $\phi_u^{(t)}$\;
    retain $B_u^{(t)}$ with \cref{eq:spring-buffer}\;
  }
\end{algorithm}

\paragraph{Inputs and outputs.}
Training consumes one user's initial adapter and a chronological stream of
query--response batches.  Each period outputs a new adapter and memory
snapshot.  Inference consumes the current adapter, current memory, and a new
query, and outputs a response.  The main experiment uses 100 LongLaMP users,
five chronological periods, Gemma-3-4B-IT, rank-8 LoRA, AdamW, two epochs,
and a 90/10 update/evaluation split within each period.

\paragraph{What transfers.}
SPRInG is the closest end-to-end prior because it preserves user ownership
through time.  Its main gap is the learning signal: difficult-to-predict text
is treated as preference drift, and updates still maximize likelihood.  Our
agent receives richer evidence.  A completed task can include tool calls,
environment outcomes, user corrections, and explicit preferences.  Those
signals should determine eligibility and reward; likelihood novelty alone
cannot tell whether an action helped the user.

\subsection{Concrete design for our policy}

The two papers suggest a clean baseline for our candidate online updates:
\begin{enumerate}
  \item Serialize user $u$'s thoughts, conversations, and successful task
  records into a chronological offline history.
  \item Train a user-owned adapter with the OPPU objective.  Save the adapter,
  optimizer state, data boundary, and base-model hash as checkpoint $c_u^{(0)}$.
  \item Collect later tasks only with the active user checkpoint.  Preserve
  states, actions, tools, outcomes, corrections, and behavior log-probabilities.
  \item Use a SPRInG-style selection stage, but score evidence by task
  validity, preference information, and provenance rather than likelihood
  novelty alone.
  \item Apply the chosen update from this note.  Activate it only after the
  shared validation and rollback gate.
\end{enumerate}

This baseline separates inheritance from novelty.  OPPU supplies the offline
per-user policy.  SPRInG supplies the chronological adapter lineage.  Our
contribution begins where supervised text prediction ends: agent evidence,
preference- or outcome-driven updates, auditable optimizer state, and safe
activation.
